\begin{abstract}
\noindent
Praca bada tezę, że Polszczyzna Kontrolowana (PK) zmniejsza złożoność obliczeniową pracy
z tekstem. W części formalnej język jest parą: zbiór zdań i relacja interpretacji.
Reguła PK jest operatorem, który usuwa zdania albo usuwa odczytania; operatory są
przemienne. Liczba odczytań zdania jest iloczynem po miejscach wyboru. Praca dowodzi
twierdzenia granicznego. Spójność zbioru wymagań bez przeczeń jest rozstrzygalna w czasie
liniowym. Reguła C5 dopuszcza 1 przeczenie w zdaniu, a 1 przeczenie wystarcza do
NP-zupełności. Część pomiarowa zawiera 5 pomiarów. Zdanie z 8 frazami przyimkowymi ma
\pkAttachEight{} odczytania, a zdanie PK ma 1 odczytanie. PK wydłuża tekst: mediana narzutu
wynosi \pkTokMedian{} \% tokenów. Sonda zadaje \pkProbeModels{} modelom językowym pytania
o sytuacje, w których odczytania zdania się różnią. Modele toru Claude odczytują zdania PK
zgodnie z kluczem w co najmniej \pkAccPkClaudeMin{} \% przypadków. Modele lokalne, które
mają od \pkLocalParamMin{} do \pkLocalParamMax{} mld parametrów włącznie, osiągają od \pkAccPkLocalMin{} \%
do \pkAccPkLocalMax{} \% włącznie. Konwencja niezgodna z nawykiem językowym działa na
model dopiero z legendą. Akty prawne mają od \pkCorpusActMin{} do \pkCorpusActMax{}
trafień sprawdzarki na 1000 słów włącznie. Część filozoficzna opisuje PK jako zestaw
kategorii obowiązkowych i rozdziela hipotezę Sapira-Whorfa dla modelu językowego na
3 tezy. Wniosek: PK zmniejsza koszt wyboru odczytania i koszt wnioskowania, a zwiększa
koszt przejścia tekstu przez model.

\smallskip
\noindent\textbf{Słowa kluczowe:} język kontrolowany; niejednoznaczność; złożoność
obliczeniowa; klauzule Horna; modele językowe; hipoteza Sapira-Whorfa.
\end{abstract}
